\section{Derivations for Gaussian smoothing selection}
\label{app:likelihood-proofs}
The following proofs record the exact fixed-reference
Gaussian calculations used in
Section~\ref{sec:gaussian-likelihood}.  They are
included to preserve the full mathematical argument
without interrupting the statistical selection narrative.

\subsection{Gaussian error-contrast likelihood}

\begin{proof}[Proof of Theorem~\ref{thm:reml}]
First use $D_dW=0$, $D_dp_\mu=\mu\mathbf1_r$,
$D_dV=U\Sigma$, and \eqref{eq:likelihood-decomposition}:
\begin{equation}\label{eq:likelihood-difference-innovation}
D_d\tau-\mu\mathbf1_r=U\Sigma\theta.
\end{equation}
Because $U$ is orthogonal and
$\operatorname{Var}(\theta)=
\sigma_u^2\Sigma^{-2}$, the covariance of
$U\Sigma\theta$ is $\sigma_u^2I_r$.
This establishes the difference-innovation interpretation.

Next multiply the observation equation by $V^\top$.
Since $V^\top W=0$ and $V^\top V=I_r$,
\begin{equation}\label{eq:contrast-observation-projection}
z=\theta+V^\top\varepsilon.
\end{equation}
Orthogonal projection of isotropic Gaussian
noise gives $V^\top\varepsilon\sim
N_r(0,\sigma_\varepsilon^2I_r)$.
The two summands are independent, so
\begin{equation}\label{eq:contrast-covariance-step}
\operatorname{Var}(z)=
\sigma_u^2\Sigma^{-2}
+\sigma_\varepsilon^2I_r
=\sigma_\varepsilon^2
\operatorname{diag}(c_1,\ldots,c_r).
\end{equation}
Consequently the $z_j$ are independent
Gaussian contrasts. The normalized negative
twice log likelihood is
\begin{equation}\label{eq:reml-logdensity}
\begin{aligned}
-2\ell_R={}&
r\log(2\pi\sigma_\varepsilon^2)
+\sum_{j=1}^r\log c_j\\
&+\frac1{\sigma_\varepsilon^2}
\sum_{j=1}^r\frac{z_j^2}{c_j}.
\end{aligned}
\end{equation}
Let $J_\lambda=\sum_jz_j^2/c_j$.
For fixed $\lambda$, differentiating
$r\log\sigma_\varepsilon^2+
J_\lambda/\sigma_\varepsilon^2$
with respect to $\sigma_\varepsilon^2$
gives $r/\sigma_\varepsilon^2
-J_\lambda/(\sigma_\varepsilon^2)^2=0$.
Since $z\ne0$, $J_\lambda>0$, and the
minimizer is
\begin{equation}\label{eq:reml-profile-step}
\widehat{\sigma}_\varepsilon^2(\lambda)
=\frac{J_\lambda}{r}.
\end{equation}
Substitution into \eqref{eq:reml-logdensity}
gives \eqref{eq:ml} plus
$r[1+\log(2\pi)]$, which does not depend on
$\lambda$.

Finally, $z$ is a full-rank orthonormal set
of contrasts orthogonal to the fixed-effect
design matrix $W$. The remaining coordinates
$W^\top(y-p_\mu)$ contain $\beta$ and are
discarded by restricted likelihood. Equivalently,
integrating the joint Gaussian density over
the unknown $\beta$ with respect to Lebesgue
measure removes those $d$ coordinates.
The resulting density is proportional to the
density of $z$, with a factor independent
of $\lambda$. This is the REML construction
for the stated model.
\end{proof}

\subsection{Quadratic and determinant identities}

\begin{proof}[Proof of Proposition~\ref{prop:det}]
Write $\tau=p_\mu+Wb+Vt$ with
$t\in\R^r$. Orthogonality separates
the PLS objective into
\begin{equation}\label{eq:orthogonal-objective}
\|W^\top(y-p_\mu)-b\|^2+
\sum_{j=1}^r\left[
(z_j-t_j)^2+\lambda\delta_jt_j^2\right].
\end{equation}
The first term is minimized by
$b=W^\top(y-p_\mu)$ and contributes
zero. For the $j$th remaining term,
differentiation gives
\begin{equation}\label{eq:orthogonal-first-order-condition}
-2(z_j-t_j)+2\lambda\delta_jt_j=0,
\quad
\widehat t_j=\frac{z_j}{1+\lambda\delta_j}.
\end{equation}
The corresponding \emph{minimized}
quadratic is
\begin{equation}\label{eq:orthogonal-minimum}
(z_j-\widehat t_j)^2+
\lambda\delta_j\widehat t_j^2
=\frac{\lambda\delta_j}{1+\lambda\delta_j}
z_j^2=\frac{z_j^2}{c_j}.
\end{equation}
Summing proves \eqref{eq:quadraticidentity}.
Furthermore, $B_d=U\Sigma^2U^\top$
has $r$ positive eigenvalues $\delta_j$,
whereas $A_\lambda$ has those
$r$ eigenvalues transformed into
$1+\lambda\delta_j$ plus $d$ unit
eigenvalues. Hence
\begin{equation}\label{eq:logdet-spectral-proof}
\begin{aligned}
&\log\det A_\lambda-r\log\lambda-\log\det B_d\\
&=\sum_{j=1}^r\log\!\left(
\frac{1+\lambda\delta_j}{\lambda\delta_j}
\right)
=\sum_{j=1}^r\log c_j,
\end{aligned}
\end{equation}
proving \eqref{eq:detidentity}.
Both $A_\lambda$ and $B_d$ are SPD
and banded, so their log determinants
follow from the log diagonals of
their banded Cholesky factors.
\end{proof}
